\documentclass[10pt,a4paper]{amsart}
\usepackage{amsmath,amssymb,amsthm,graphicx}
% Retained preamble for the Stacks Project chapter snapshot in this directory.
% The chapter's own preamble.tex is a copy of the shared Stacks Project preamble; the block
% below is sliced MECHANICALLY, line for line, from it. Nothing was added, removed or
% redefined; the wrapper class, the xr-hyper/hyperref wiring and the externaldocument list
% are omitted because the assembler supplies the document class.
%
% lines taken from preamble.tex: 3, 5, 9, 11, 12, 156, 157, 158, ... (51 lines total)
%
% package coverage restored from the chapter's own preamble (2026-10-13): xy, xy-2cell

\usepackage[all]{xy}
\xyoption{2cell}
\UseAllTwocells

\newenvironment{reference}{\comment}{\endcomment}
\newenvironment{slogan}{\comment}{\endcomment}
\newenvironment{history}{\comment}{\endcomment}
\theoremstyle{plain}
\newtheorem{theorem}[subsection]{Theorem}
\newtheorem{proposition}[subsection]{Proposition}
\newtheorem{lemma}[subsection]{Lemma}
\theoremstyle{definition}
\newtheorem{definition}[subsection]{Definition}
\newtheorem{example}[subsection]{Example}
\newtheorem{exercise}[subsection]{Exercise}
\newtheorem{situation}[subsection]{Situation}
\theoremstyle{remark}
\newtheorem{remark}[subsection]{Remark}
\newtheorem{remarks}[subsection]{Remarks}
\numberwithin{equation}{subsection}
\def\lim{\mathop{\mathrm{lim}}\nolimits}
\def\colim{\mathop{\mathrm{colim}}\nolimits}
\def\Spec{\mathop{\mathrm{Spec}}}
\def\Hom{\mathop{\mathrm{Hom}}\nolimits}
\def\Ext{\mathop{\mathrm{Ext}}\nolimits}
\def\SheafHom{\mathop{\mathcal{H}\!\mathit{om}}\nolimits}
\def\SheafExt{\mathop{\mathcal{E}\!\mathit{xt}}\nolimits}
\def\Sch{\mathit{Sch}}
\def\Mor{\mathop{\mathrm{Mor}}\nolimits}
\def\Ob{\mathop{\mathrm{Ob}}\nolimits}
\def\Sh{\mathop{\mathit{Sh}}\nolimits}
\def\NL{\mathop{N\!L}\nolimits}
\def\CH{\mathop{\mathrm{CH}}\nolimits}
\def\proetale{{pro\text{-}\acute{e}tale}}
\def\etale{{\acute{e}tale}}
\def\QCoh{\mathit{QCoh}}
\def\Ker{\mathop{\mathrm{Ker}}}
\def\Im{\mathop{\mathrm{Im}}}
\def\Coker{\mathop{\mathrm{Coker}}}
\def\Coim{\mathop{\mathrm{Coim}}}
\DeclareMathSymbol{\boxtimes}{\mathbin}{AMSa}{"02}
\def\QCohstack{\mathcal{QC}\!\mathit{oh}}
\def\Cohstack{\mathcal{C}\!\mathit{oh}}
\def\Spacesstack{\mathcal{S}\!\mathit{paces}}
\def\Quotfunctor{\mathrm{Quot}}
\def\Hilbfunctor{\mathrm{Hilb}}
\def\Curvesstack{\mathcal{C}\!\mathit{urves}}
\def\Polarizedstack{\mathcal{P}\!\mathit{olarized}}
\def\Complexesstack{\mathcal{C}\!\mathit{omplexes}}
\def\Pic{\mathop{\mathrm{Pic}}\nolimits}
\def\Picardstack{\mathcal{P}\!\mathit{ic}}
\def\Picardfunctor{\mathrm{Pic}}
\def\Deformationcategory{\mathcal{D}\!\mathit{ef}}
\usepackage{cleveref}
\crefname{originallemma}{original lemma}{original lemmas}
\Crefname{originallemma}{Original lemma}{Original lemmas}
\makeatletter
\newlabel{example-rings}{{0DY1}{source}{}{sourcepointer-1}{}}
\newlabel{example-rings@cref}{{[originallemma][0][]0DY1}{[1][0][]source}{}{}{}}
\makeatother
\crefname{originallemma}{original lemma}{original lemmas}
\Crefname{originallemma}{Original lemma}{Original lemmas}
\makeatletter
\newlabel{example-schemes}{{0DY7}{source}{}{sourcepointer-2}{}}
\newlabel{example-schemes@cref}{{[originallemma][0][]0DY7}{[1][0][]source}{}{}{}}
\makeatother
\crefname{originallemma}{original lemma}{original lemmas}
\Crefname{originallemma}{Original lemma}{Original lemmas}
\makeatletter
\newlabel{lemma-affine}{{0DYX}{source}{}{sourcepointer-3}{}}
\newlabel{lemma-affine@cref}{{[originallemma][0][]0DYX}{[1][0][]source}{}{}{}}
\makeatother
\crefname{originallemma}{original lemma}{original lemmas}
\Crefname{originallemma}{Original lemma}{Original lemmas}
\makeatletter
\newlabel{lemma-open}{{0DYW}{source}{}{sourcepointer-4}{}}
\newlabel{lemma-open@cref}{{[originallemma][0][]0DYW}{[1][0][]source}{}{}{}}
\makeatother
\crefname{originallemma}{original lemma}{original lemmas}
\Crefname{originallemma}{Original lemma}{Original lemmas}
\makeatletter
\newlabel{schemes-lemma-glue}{{01JB}{source}{}{sourcepointer-5}{}}
\newlabel{schemes-lemma-glue@cref}{{[originallemma][0][]01JB}{[1][0][]source}{}{}{}}
\makeatother
\crefname{originallemma}{original lemma}{original lemmas}
\Crefname{originallemma}{Original lemma}{Original lemmas}
\makeatletter
\newlabel{schemes-section-glueing-schemes}{{01JA}{source}{}{sourcepointer-6}{}}
\newlabel{schemes-section-glueing-schemes@cref}{{[originallemma][0][]01JA}{[1][0][]source}{}{}{}}
\makeatother
\crefname{originallemma}{original lemma}{original lemmas}
\Crefname{originallemma}{Original lemma}{Original lemmas}
\makeatletter
\newlabel{situation-glueing}{{0DYY}{source}{}{sourcepointer-7}{}}
\newlabel{situation-glueing@cref}{{[originallemma][0][]0DYY}{[1][0][]source}{}{}{}}
\makeatother
\title{Stacks, Examples of Deformations: lemma-glueing}
\author{The Stacks Project Authors}
\date{stacks-project:examples-defos.tex:local-snapshot-20261012; GNU FDL 1.2 or later}
\begin{document}
\maketitle
\noindent\textit{Editorial note.} Statement span is lines 1650-1662 of sources/stacks-examples-defos/pinned-original-source.tex; proof span is lines 1664-1698, the complete author proof (33 body lines).\par\medskip
\noindent\textit{Editorial note.} Navigation stubs: 7 cited-but-not-included labels are supplied as non-load-bearing pointers with their pinned tags where the snapshot tag index has them.\par\medskip
\noindent\textit{Editorial note.} extra\_packages is empty on purpose: the retained preamble already loads xy with options and a second declaration clashes in the assembler wrapper.\par\medskip
\noindent\textit{Editorial note.} Text is retained verbatim, including the author's own typos and placeholders.\par\medskip
\noindent\textit{Editorial note.} Queue index 17 of queue\_examples-defos.json (pb=1425, stmt 1650-1662, proof 1664-1698).\par\medskip
\noindent\textit{Allowed standard prerequisites (editorial).}\begin{itemize}
\item Explicit prerequisite boundary. The cited inputs are prerequisites; the argument in the retained span is not replaced by any cited result.
\end{itemize}\medskip
\ifdefined\hypertarget\hypertarget{sourcepointer-1}{}\fi
\noindent\textit{Source reference.} Original lemma 0DY1: sources/stacks-examples-defos/pinned-original-source.tex, source lines 1031--1040. Cited by the retained statement or proof but not part of the collected spans; supplied only as a navigation pointer. The argument does not depend on its content.\par\medskip
\ifdefined\hypertarget\hypertarget{sourcepointer-2}{}\fi
\noindent\textit{Source reference.} Original lemma 0DY7: sources/stacks-examples-defos/pinned-original-source.tex, source lines 1425--1440. Cited by the retained statement or proof but not part of the collected spans; supplied only as a navigation pointer. The argument does not depend on its content.\par\medskip
\ifdefined\hypertarget\hypertarget{sourcepointer-3}{}\fi
\noindent\textit{Source reference.} Original lemma 0DYX: sources/stacks-examples-defos/pinned-original-source.tex, source lines 1600--1608. Cited by the retained statement or proof but not part of the collected spans; supplied only as a navigation pointer. The argument does not depend on its content.\par\medskip
\ifdefined\hypertarget\hypertarget{sourcepointer-4}{}\fi
\noindent\textit{Source reference.} Original lemma 0DYW: sources/stacks-examples-defos/pinned-original-source.tex, source lines 1584--1592. Cited by the retained statement or proof but not part of the collected spans; supplied only as a navigation pointer. The argument does not depend on its content.\par\medskip
\ifdefined\hypertarget\hypertarget{sourcepointer-5}{}\fi
\noindent\textit{Source reference.} Original lemma 01JB: sources/stacks-examples-defos/pinned-original-source.tex, source lines 1698--1698. Cited by the retained statement or proof but not part of the collected spans; supplied only as a navigation pointer. The argument does not depend on its content.\par\medskip
\ifdefined\hypertarget\hypertarget{sourcepointer-6}{}\fi
\noindent\textit{Source reference.} Original lemma 01JA: sources/stacks-examples-defos/pinned-original-source.tex, source lines 1698--1698. Cited by the retained statement or proof but not part of the collected spans; supplied only as a navigation pointer. The argument does not depend on its content.\par\medskip
\ifdefined\hypertarget\hypertarget{sourcepointer-7}{}\fi
\noindent\textit{Source reference.} Original lemma 0DYY: sources/stacks-examples-defos/pinned-original-source.tex, source lines 1638--1648. Cited by the retained statement or proof but not part of the collected spans; supplied only as a navigation pointer. The argument does not depend on its content.\par\medskip
\begin{lemma}
\label{lemma-glueing}
In Situation \ref{situation-glueing}
there is an equivalence
$$
\Deformationcategory_X =
\Deformationcategory_{P_1}
\times_{\Deformationcategory_{P_{12}}}
\Deformationcategory_{P_2}
$$
of deformation categories, see Examples \ref{example-schemes} and
\ref{example-rings}.
\end{lemma}

\begin{proof}
It suffices to show that the functors of Lemma \ref{lemma-open}
define an equivalence
$$
\Deformationcategory_X \longrightarrow
\Deformationcategory_{U_1}
\times_{\Deformationcategory_{U_{12}}}
\Deformationcategory_{U_2}
$$
because then we can apply Lemma \ref{lemma-affine} to translate into rings.
To do this we construct a quasi-inverse. Denote
$F_i : \Deformationcategory_{U_i} \to \Deformationcategory_{U_{12}}$
the functor of Lemma \ref{lemma-open}.
An object of the RHS is given by an $A$ in $\mathcal{C}_\Lambda$,
objects $(A, V_1) \to (k, U_1)$ and $(A, V_2) \to (k, U_2)$, and
a morphism
$$
g : F_1(A, V_1) \to F_2(A, V_2)
$$
Now $F_i(A, V_i) = (A, V_{i, 3 - i})$ where $V_{i, 3 - i} \subset V_i$
is the open subscheme whose base change to $k$ is $U_{12} \subset U_i$.
The morphism $g$ defines an isomorphism
$V_{1, 2} \to V_{2, 1}$ of schemes over $A$ compatible
with $\text{id} : U_{12} \to U_{12}$ over $k$.
Thus $(\{1, 2\}, V_i, V_{i, 3 - i}, g, g^{-1})$ is a glueing
data as in Schemes, Section \ref{schemes-section-glueing-schemes}.
Let $Y$ be the glueing, see Schemes, Lemma \ref{schemes-lemma-glue}.
Then $Y$ is a scheme over $A$ and the
compatibilities mentioned above show that
there is a canonical isomorphism
$Y \times_{\Spec(A)} \Spec(k) = X$.
Thus $(A, Y) \to (k, X)$ is an object of $\Deformationcategory_X$.
We omit the verification that this construction is a functor
and is quasi-inverse to the given one.
\end{proof}

\end{document}
